%%% FIELD def-contrast-class-construction
\(\mathcal C_K=\{c\in\mathbb R^K: \sum_{a\in\mathcal A_K}c_a=0,\ \sum_{a\in\mathcal A_K}|c_a|=2\}\).
%%% FIELD statement-response-gauge
For every \(c\in\mathcal C_K\), the primal program defining \(V_c\) is feasible for every \(\xi\in\mathbb R^{S_c}\), and
\[
V_c(\xi)=\max\{\ell^\top\xi:|\ell^\top b_t|\le\theta_t^2\text{ for every }t\in\mathcal T_c\}.
\]
The dual feasible set is compact on \(\mathbf 1^\perp\), \(V_c\) is finite, convex, even, continuous, and positively homogeneous, and \(0<\kappa(c)\le1\). For rational \(c\), primal and dual optimizers and \(\kappa(c)\) are exactly rational and terminate under finite linear programming. These conclusions include the support-two case.
%%% FIELD proof-response-gauge
Let \(d=|S_c|\).  The zero-sum and \(\ell^1\)-normalization assumptions imply
\[
\sum_{a:c_a>0}c_a=1,
\qquad
\sum_{a:c_a<0}c_a=-1,
\tag{1}
\]
so \(c\) has at least one positive and one negative coordinate and hence \(d\ge2\).  Regard the vectors \(b_t\) as the columns of a \(d\)-by-\(2^d\) matrix \(A\).  If \(e_a\) is an active-coordinate unit vector, then
\[
e_a=b_{e_a}+b_{\mathbf 1},
\qquad
-e_a=b_{\mathbf 1-e_a}+b_{\mathbf 0}.
\tag{2}
\]
For \(\xi_a^+=\max\{\xi_a,0\}\) and \(\xi_a^-=\max\{-\xi_a,0\}\), multiplying the first identity by \(\xi_a^+\), the second by \(\xi_a^-\), and summing over \(a\) gives a nonnegative response-type representation of \(\xi\).  Thus the primal program is feasible.  Its cost is finite because \(\mathcal T_c\) is finite.

Put \(w_t=\theta_t^2\).  The primal is
\[
p(\xi)=\min\{w^\top\lambda:A\lambda=\xi,\ \lambda\ge0\}.
\]
Weak duality gives
\[
\sup\{\ell^\top\xi:A^\top\ell\le w\}\le p(\xi).
\tag{3}
\]
We prove the reverse inequality rather than assume linear-programming duality.  The set
\[
\mathcal P=\{(A\lambda,w^\top\lambda+s):\lambda\ge0,\ s\ge0\}
\]
is a closed polyhedral cone.  For \(\varepsilon>0\), the point \((\xi,p(\xi)-\varepsilon)\) is not in \(\mathcal P\), so finite-dimensional strict separation gives a separating vector.  Because \(\mathcal P\) contains the vertical ray above \((\xi,p(\xi))\), the last coordinate of this vector is nonpositive; it cannot be zero because the separated points have the same first coordinate.  Normalize it to \(-1\).  The resulting first coordinate \(\ell_\varepsilon\) obeys \(A^\top\ell_\varepsilon\le w\) and \(\ell_\varepsilon^\top\xi\ge p(\xi)-\varepsilon\).  Letting \(\varepsilon\downarrow0\) in (3) proves equality of the primal and dual values.

The complement map \(t\mapsto\mathbf 1-t\) sends \(b_t\) to \(-b_t\) and \(\theta_t\) to \(-\theta_t\).  Pairing complementary dual constraints therefore turns \(A^\top\ell\le w\) into
\[
|\ell^\top b_t|\le\theta_t^2
\quad(t\in\mathcal T_c).
\tag{4}
\]
The all-zero and all-one types have zero cost by the zero-sum assumption.  Their constraints imply \(\ell^\top\mathbf 1=0\).  The singleton constraint then yields
\[
|\ell_a|=|\ell^\top(e_a-\mathbf 1/2)|\le c_a^2.
\tag{5}
\]
Thus the dual feasible set is a closed subset of the bounded box \(\prod_{a\in S_c}[-c_a^2,c_a^2]\), contained in \(\mathbf1^\perp\), and is compact.  It is nonempty because \(\ell=0\) is feasible.  The dual formula identifies \(V_c\) as the support function of a nonempty compact centrally symmetric convex set.  It follows directly from the support-function formula that \(V_c\) is finite, convex, even, continuous, and positively homogeneous.

To bound \(\kappa(c)\), define
\[
m_c=\min\{|\theta_t|:t\in\mathcal T_c,\ \theta_t\ne0\}.
\]
This is positive because the alphabet is finite and, by (1), the type \(t^+\) with \(t_a^+=\mathbf1\{c_a>0\}\) has \(\theta_{t^+}=1\).  Every feasible representation of \(v\) satisfies
\[
1=c^\top v=\sum_t\lambda_tc^\top b_t=\sum_t\lambda_t\theta_t.
\]
Consequently,
\[
\sum_t\lambda_t\theta_t^2
\ge m_c\sum_t\lambda_t|\theta_t|
\ge m_c\left|\sum_t\lambda_t\theta_t\right|=m_c.
\tag{6}
\]
On the other hand, \(b_{t^+}=s_c/2=v\) and \(\theta_{t^+}=1\), so the one-atom representation \(\lambda_{t^+}=1\) has cost one.  Taking the infimum in (6) proves \(0<m_c\le\kappa(c)\le1\).

Finally, if \(c\) is rational, then \(A\), \(v\), and \(w\) are rational.  Starting from any optimal feasible \(\lambda\), if its positive-support columns are dependent, move in a nonzero rational null direction and choose the sign that does not increase cost until one positive coordinate reaches zero.  Repetition produces an optimal basic feasible solution.  There are finitely many column supports; exact Gaussian elimination enumerates all their rational vertices and compares their rational costs.  The compact dual polytope is rational as well, and successive rational linear programs select its lexicographically least optimal point.  Hence rational primal and dual optimizers and the rational value \(\kappa(c)\) are found after finitely many exact operations.  No step used \(d\ge3\), so the proof includes \(d=2\).
%%% FIELD statement-universal-upper
For every fixed \(K\ge2\) and every \(c\in\mathcal C_K\), the independent allocation in \(\mathrm{def:independent\text{-}qstar\text{-}design}\) and the total estimator \(\delta_n^\star\) satisfy
\[
\sup_{y\in\{0,1\}^{n\times K}}\mathcal R_n(\mathcal D_n^\star,\delta_n^\star;y)
\le n^{-1}-\Lambda_cn^{-4/3}+o(n^{-4/3}).
\]
Equivalently, \(\liminf_{n\to\infty}n^{4/3}d_n(c)\ge\Lambda_c\).  The design assigns only active arms, so the assertion also covers contrasts with two active arms and any number of inactive arms.
%%% FIELD proof-universal-upper
Fix a schedule \(y\), and put
\[
t_i=(Y_i(a))_{a\in S_c},\qquad
\theta_i=c^\top t_i,\qquad
u_i=(g^\star)^\top b_{t_i}.
\]
The zero-sum assumption gives \(c^\top t_i=c^\top b_{t_i}\), and hence \(\tau_c(y)=n^{-1}\sum_i\theta_i\).  Equation (1) in the gauge proof gives \(\sum_{a\in S_c}q_a^\star=1\).  Thus \(\mathcal D_n^\star\) is a probability law, and it puts probability zero on every inactive arm.  Under this product design, the score definitions give
\[
\mathbb ET_i=\theta_i,
\qquad
\mathbb EU_i=u_i,
\tag{7}
\]
and, because \((Y_i^{\mathrm{obs}}-1/2)^2=1/4\),
\[
\mathbb ET_i^2=\sum_{a\in S_c}\frac{c_a^2}{4q_a^\star}=1,
\qquad
\mathbb E(T_iU_i)=\sum_{a\in S_c}\frac{c_ag_a^\star}{4q_a^\star}
=(g^\star)^\top v=1.
\tag{8}
\]
Here the last equality uses \((\ell^\star)^\top v=\kappa(c)\) from primal-dual equality and \(g^\star=\ell^\star/\kappa(c)\).  The dual constraints further imply
\[
|u_i|\le\theta_i^2/\kappa(c).
\tag{9}
\]
All variables in (7)--(9) are bounded by constants depending only on the fixed contrast.  Positivity of every active \(q_a^\star=|c_a|/2\) justifies each division.

Let \(\psi\) and \(h=-2\psi'/\psi\) be supplied by \(\mathrm{lem:unit\text{-}airy\text{-}feedback\text{-}properties}\).  Its Riccati identity is equivalent in the weak sense to
\[
-4\psi''+|x|\psi=C_A\psi.
\]
For compactly supported smooth \(f\), integration by parts gives the ground-state identity
\[
\int_{\mathbb R}\{4|f'|^2+(|x|-C_A)f^2\}\,dx
=4\int_{\mathbb R}\psi^2\left|\left(\frac f\psi\right)'\right|^2dx\ge0.
\tag{10}
\]
Approximation extends (10) to the variational domain, while equality holds for \(f=\psi\).  Thus \(\psi\) is the unit-scale minimizer in the definition of \(C_A\).  Scaling the variational problem yields
\[
\phi_\kappa(x)=\kappa(c)^{1/6}\psi(\kappa(c)^{1/3}x),
\qquad
h_\kappa(x)=\kappa(c)^{1/3}h(\kappa(c)^{1/3}x).
\]
Therefore
\[
h_\kappa(x)^2-2h_\kappa'(x)
=\kappa(c)|x|-C_A\kappa(c)^{2/3}
=\kappa(c)|x|-\Lambda_c.
\tag{11}
\]
The cited unit-scale bounds and this scaling also give, for a finite constant \(M_c\),
\[
|h_\kappa'(x)|\le\frac{M_c}{1+\sqrt{|x|}},
\tag{12}
\]
and show that \(h_\kappa^2\) and \(h_\kappa'\) are globally Lipschitz.

Set
\[
X_i=T_i-\theta_i,\quad W_i=U_i-u_i,\quad
A_n=\sum_iX_i,\quad x_n=n^{-2/3}\sum_i u_i,\quad
S_n=n^{-2/3}\sum_iU_i.
\]
Independence across units is part of \(\mathcal D_n^\star\), so (7)--(8) imply
\[
\mathbb EA_n^2=\sum_i(1-\theta_i^2).
\tag{13}
\]
For \(S_{-i}=x_n+n^{-2/3}\sum_{j\ne i}W_j\), Taylor's formula with Lipschitz derivative, independence of unit \(i\) from \(S_{-i}\), and (8) give
\[
\mathbb E\{X_ih_\kappa(S_n)\}
=n^{-2/3}(1-\theta_iu_i)\mathbb Eh_\kappa'(S_{-i})+O_c(n^{-4/3}).
\tag{14}
\]
Since the bounded independent centered variables satisfy
\[
\mathbb E|S_{-i}-x_n|
\le\{\mathbb E(S_{-i}-x_n)^2\}^{1/2}=O_c(n^{-1/6}),
\]
the Lipschitz property of \(h_\kappa'\), followed by summation of (14), yields uniformly in \(y\)
\[
\mathbb E\{A_nh_\kappa(S_n)\}
=n^{1/3}h_\kappa'(x_n)
-n^{-2/3}h_\kappa'(x_n)\sum_i\theta_iu_i
+O_c(n^{1/6}).
\tag{15}
\]
Similarly, the global Lipschitz property of \(h_\kappa^2\) gives
\[
\mathbb Eh_\kappa(S_n)^2=h_\kappa(x_n)^2+O_c(n^{-1/6}).
\tag{16}
\]

Before projection, expanding squared error and applying (13)--(16) gives
\[
\begin{split}
\mathbb E\left[n^{-1}\sum_iT_i-n^{-2/3}h_\kappa(S_n)-\tau_c(y)\right]^2
={}&\frac1n-\frac1{n^2}\sum_i\theta_i^2
+n^{-4/3}\{h_\kappa(x_n)^2-2h_\kappa'(x_n)\}\\
&+2n^{-7/3}h_\kappa'(x_n)\sum_i\theta_iu_i+O_c(n^{-3/2}).
\end{split}
\tag{17}
\]
By (9),
\[
\sum_i\theta_i^2\ge\kappa(c)\left|\sum_i u_i\right|
=\kappa(c)n^{2/3}|x_n|.
\tag{18}
\]
Using (9), (12), (18), and the deterministic bound \(|x_n|=O_c(n^{1/3})\), we obtain
\[
-n^{-2}\sum_i\theta_i^2
+2n^{-7/3}h_\kappa'(x_n)\sum_i\theta_iu_i
\le-\kappa(c)|x_n|n^{-4/3}+O_c(n^{-3/2}).
\tag{19}
\]
Indeed, after the leading negative term is bounded with (18), the possible loss is at most
\(C_cn^{-5/3}|x_n|/(1+\sqrt{|x_n|})=O_c(n^{-3/2})\).

Substituting (11) and (19) into (17) proves
\[
\sup_y\mathbb E\left[n^{-1}\sum_iT_i-n^{-2/3}h_\kappa(S_n)-\tau_c(y)\right]^2
\le n^{-1}-\Lambda_cn^{-4/3}+O_c(n^{-3/2}).
\]
Since \(\tau_c(y)\in[-1,1]\), metric projection onto \([-1,1]\) cannot increase squared error.  This proves the displayed risk bound for \(\delta_n^\star\).  The minimax definition then gives \(R_n(c)\le\sup_y\mathcal R_n(\mathcal D_n^\star,\delta_n^\star;y)\), which is equivalent to the claimed lower limit for \(n^{4/3}d_n(c)\).  The argument used only active labels and therefore remains valid when \(d=2\) and when \(K-d\) inactive labels are present.
%%% FIELD statement-all-count
For every fixed \(K\ge2\) and \(c\in\mathcal C_K\), there is a sequence of permutation-invariant priors \(\Pi_n\) on the finite set \(\{0,1\}^{n\times K}\), chosen without reference to the assignment law, such that
\[
\inf_{r:\,\sum_ar_a=n}\operatorname{BR}_{\Pi_n}(r)
\ge n^{-1}-\Lambda_cn^{-4/3}-o(n^{-4/3}).
\]
The remainder is uniform over all count vectors, including allocation to inactive arms and every face on which one or more active arms have zero count.  Conditional on every fixed labeled assignment, the posterior regression is the exact regression of the random finite-population target \(\tau_c(y)\); the inequality remains valid after averaging under every possibly dependent outcome-independent design \(\mathcal D_n\) and against every randomized full-label estimator \(\widehat\tau\).  When \(|S_c|=2\), the construction has one transverse coordinate and the same proof covers all inactive-arm and missing-active-arm faces.
%%% FIELD proof-all-count
Write \(d=|S_c|\).  As shown in (1) of the gauge proof, the positive and negative coefficients sum to \(1\) and \(-1\), respectively, so \(d\ge2\).  Relabel the active coordinates as \(a_0,a_1,\ldots,a_{d-1}\), and set
\[
d_j=e_{a_j}-\frac{c_{a_j}}{c_{a_0}}e_{a_0},
\qquad j=1,\ldots,d-1.
\tag{20}
\]
The denominator in (20) is nonzero because \(a_0\in S_c\).  Each \(d_j\) lies in \(\ker(c^\top)\), and the \(d_j\)'s form a basis of that kernel.  Moreover, \(B=[v,d_1,\ldots,d_{d-1}]\) is invertible: applying \(c^\top\) to a null linear combination first makes its coefficient of \(v\) zero because \(c^\top v=1\), and independence of the \(d_j\)'s then kills the remaining coefficients.  The first row of \(B^{-1}\) is \(c^\top\).  For \(d=2\), (20) simply supplies the single transverse direction required below.

We first construct a proper response tube.  The response identities (2) provide a finite nonnegative response-type representation of every vector.  Choose an optimal representation \(\lambda^{v,+}\) of \(v\), and choose any such representation \(\lambda^{j,+}\) of each \(d_j\).  Complementing response types gives representations \(\lambda^{v,-}\) and \(\lambda^{j,-}\) of the negative directions with the same total masses and costs.  Write these total masses as \(L_v,L_j\), and write
\[
k_v=\sum_t\lambda_t^{v,+}\theta_t^2=\kappa(c),
\qquad
k_j=\sum_t\lambda_t^{j,+}\theta_t^2.
\]
Let \(p^0\) put mass \(1/2\) on each of the all-zero and all-one response types.  Fix \(T<\infty\) and \(\varepsilon>0\), put \(\alpha=n^{-1/3}\), and, for \(x\in[-T,T]\) and \(h_j\in[-\varepsilon,\varepsilon]\), define
\[
\begin{split}
p_t(x,h)=p_t^0+\alpha\bigg[&|x|\lambda_t^{v,\operatorname{sgn}(x)}
+\sum_j|h_j|\lambda_t^{j,\operatorname{sgn}(h_j)}\\
&-\left(L_v|x|+\sum_jL_j|h_j|\right)p_t^0\bigg].
\end{split}
\tag{21}
\]
At a zero coordinate the associated term is zero, so its selected sign is immaterial.  The bracket has total mass zero.  If
\[
\alpha\left(L_vT+\varepsilon\sum_jL_j\right)\le\frac12,
\tag{22}
\]
then every noncentral probability is nonnegative and each central probability is at least \(1/4\); hence (21) is a probability vector.  If \(t\) is drawn from this vector and \(\zeta=xv+\sum_jh_jd_j\), the representation identities give exactly
\[
\mathbb Eb_t=\alpha\zeta,
\qquad
\mathbb E\theta_t=\alpha x,
\qquad
\operatorname{Var}(\theta_t)
=\alpha\left\{\kappa(c)|x|+\sum_jk_j|h_j|\right\}-\alpha^2x^2.
\tag{23}
\]

Choose a compactly supported continuous piecewise-polynomial \(f\), positive inside its support and zero at its endpoints, and put
\[
w(x)=\frac{f(x)^2}{\int f^2},
\qquad
J_x=\int\frac{w'(x)^2}{w(x)}\,dx
=4\frac{\int f'(x)^2\,dx}{\int f(x)^2\,dx}.
\]
Scaling \(x\mapsto\kappa(c)^{1/3}x\) in the variational definition of \(C_A\) shows that the infimum of \(J_x+\kappa(c)\mathbb E_w|x|\) is \(\Lambda_c\).  Compact truncation, endpoint tapering, and rational piecewise-polynomial interpolation in the weighted Sobolev norm therefore give, for every \(\epsilon>0\), an \(f\) with rational knots and coefficients such that
\[
J_x+\kappa(c)\mathbb E_w|x|\le\Lambda_c+\epsilon.
\tag{24}
\]
The operations just listed preserve convergence of \(\int f^2\), \(\int f'^2\), and \(\int|x|f^2\); after truncation these are ordinary integrals on a compact interval.  This proves (24) directly from the defining infimum, rather than assuming existence of an exact compact minimizer.

Independently for every transverse coordinate use
\[
\rho_\varepsilon(h)=\frac{15}{16\varepsilon}
\left(1-\frac{h^2}{\varepsilon^2}\right)^2
\mathbf1\{|h|\le\varepsilon\}.
\]
It vanishes quadratically at both endpoints, and direct integration gives
\[
J_h=\int\frac{\rho_\varepsilon'(h)^2}{\rho_\varepsilon(h)}\,dh
=\frac{10}{\varepsilon^2}.
\tag{25}
\]
Draw \((x,h)\) from \(w\prod_j\rho_\varepsilon\); conditionally draw the \(n\) active response vectors independently from (21); and draw all inactive-arm outcomes as independent fair bits, independently of everything else.  Integrating out \((x,h)\) defines a prior \(\Pi_n\) on the finite schedule space.  It is invariant under permutations of unit labels and is chosen without reference to the assignment law.

Fix any labeled assignment, put \(q_a=r_a/n\) on active arms, and let \(Q=\sum_{a\in S_c}q_a\le1\).  Conditional on \((x,h)\), define
\[
\mu=\mathbb E\theta_t=\alpha x,\quad
\mu_a=\mathbb Eb_{t,a},\quad
V_a=\operatorname{Var}(b_{t,a})=\frac14-\mu_a^2,
\]
\[
C_a=\operatorname{Cov}(\theta_t,b_{t,a}),
\qquad
\sigma^2=\operatorname{Var}(\theta_t).
\]
Compactness of the parameter box in (21) and its finite, cellwise-affine formula give constants \(L,\Gamma<\infty\), depending only on the fixed tube, such that
\[
|\mu_a|\le\alpha L,
\qquad
|C_a|\le\alpha\Gamma.
\tag{26}
\]
For the covariance bound, note that \(\mathbb E_{p^0}(\theta_tb_{t,a})=0\), and maximize the finitely many affine coefficients in (21).  For all sufficiently large \(n\), \(4\alpha^2L^2\le1/2\), so
\[
\frac18\le V_a\le\frac14,
\qquad
V_a^{-1}\le4(1+u),
\qquad
u=8\alpha^2L^2.
\tag{27}
\]
These bounds justify every division by \(V_a\).

The observed centered outcome on an active arm takes only the values \(\pm1/2\).  Every function of such a binary variable is affine; matching its mean and covariance therefore gives the exact conditional regression
\[
\mathbb E(\theta_t\mid x,h,b_{t,a})
=\mu+\frac{C_a}{V_a}(b_{t,a}-\mu_a).
\]
If \(O\) denotes all labels and observed outcomes, conditional independence across units yields
\[
m_{x,h}(O):=\mathbb E\{\tau_c(y)\mid x,h,O\}
=\mu+\frac1n\sum_{i:Z_i\in S_c}
\frac{C_{Z_i}}{V_{Z_i}}(b_{t_i,Z_i}-\mu_{Z_i}).
\tag{28}
\]
Inactive observations do not enter (28) because their prior law is independent of \((x,h)\) and of the active response types.  Applying the conditional-variance identity first given \((x,h,O)\), and then given \(O\), gives the exact target decomposition
\[
\operatorname{BR}_{\Pi_n}(r)
=\frac1n\mathbb E\left(\sigma^2-\sum_aq_a\frac{C_a^2}{V_a}\right)
+\mathbb E\operatorname{Var}\{m_{x,h}(O)\mid O\}.
\tag{29}
\]
This is the posterior variance of the random finite-population target itself, not merely the variance of its conditional mean.

We next derive a directional Bayes inequality.  Let \(D\) be a fixed directional derivative in \((x,h)\)-coordinates satisfying
\[
D\mu=\alpha,
\qquad
D\mu_a=\alpha z_a.
\]
Let \(S_D\) be the directional score of the joint parameter-observation density.  The endpoint-vanishing priors remove the outer boundary terms.  Formula (21), and hence (28), is continuous across every sign cell and has a bounded weak derivative on each cell; internal face terms cancel in opposite pairs.  The same cancellation holds at spline knots.  Applying one-dimensional integration by parts successively in the coordinates therefore gives
\[
\mathbb E\left[(m_{x,h}(O)-\mathbb E\{m_{x,h}(O)\mid O\})S_D\right]
=-\mathbb E(Dm_{x,h}(O)).
\tag{30}
\]
Indeed, conditional on \(O\), the full joint score integrates to zero, including at the quadratic-zero endpoints.  Differentiating (28), the derivative of \(C_a/V_a\) is multiplied by \(b_{t_i,a}-\mu_a\), whose conditional mean is zero.  Hence
\[
\mathbb E(Dm_{x,h}(O))
=\alpha\left[1-\sum_aq_a\mathbb E\left(\frac{z_aC_a}{V_a}\right)\right].
\tag{31}
\]
Prior and likelihood scores are orthogonal, and distinct-unit likelihood scores are conditionally orthogonal.  A Bernoulli variable with mean \(1/2+\mu_a\) has directional information \((D\mu_a)^2/V_a\).  It follows that
\[
\mathbb ES_D^2
=J_D+n\alpha^2\sum_aq_a\mathbb E(z_a^2/V_a).
\tag{32}
\]
Cauchy--Schwarz applied to (30)--(32) proves
\[
\mathbb E\operatorname{Var}\{m_{x,h}(O)\mid O\}
\ge
\frac{\alpha^2\left[1-\sum_aq_a\mathbb E(z_aC_a/V_a)\right]^2}
{J_D+n\alpha^2\sum_aq_a\mathbb E(z_a^2/V_a)}.
\tag{33}
\]
Thus (33), including its boundary behavior, has been derived for the piecewise-polynomial tube.

It remains to make (29) and (33) uniform over all allocation faces.  Put \(\delta=n^{-1/8}\).  Suppose first that
\[
\max_{a\in\mathcal A_K}|q_a-q_a^\star|\le\delta.
\tag{34}
\]
Use \(D=\partial_x\), so \(z=v\) and \(J_D=J_x\).  The exact covariance identity
\[
\sum_ac_aC_a
=\operatorname{Cov}\left(\theta_t,\sum_ac_ab_{t,a}\right)=\sigma^2
\tag{35}
\]
and \(4q_a^\star v_a=c_a\), together with (26)--(27), imply
\[
\left|\sum_aq_a\mathbb E(v_aC_a/V_a)-\mathbb E\sigma^2\right|
\le4d\delta\alpha\Gamma+16\alpha^3\Gamma L^2.
\tag{36}
\]
Also,
\[
\sum_aq_av_a^2/V_a\le1+u,
\qquad
\sum_aq_aC_a^2/V_a\le4(1+u)\alpha^2\Gamma^2.
\tag{37}
\]
Substitute (23), (33), (36), and (37) into (29), and use \((1-s)^2/(1+t)\ge1-2s-t\) for \(t\ge0\).  Direct collection of terms gives
\[
\operatorname{BR}_{\Pi_n}(r)
\ge\frac1n-n^{-4/3}\left[J_x+\kappa(c)\mathbb E|x|
+\sum_jk_j\mathbb E|h_j|+e_{\mathrm{loc}}(n)\right],
\tag{38}
\]
where, for example,
\[
e_{\mathrm{loc}}(n)=8\alpha L^2+4(1+u)\alpha\Gamma^2
+8d\delta\Gamma+32\alpha^2\Gamma L^2=o(1).
\tag{39}
\]

Next suppose every active \(q_a\) is positive and (34) fails.  Define
\[
H=\sum_{a\in S_c}\frac{c_a^2}{q_a},
\qquad
z_a=\frac{c_a}{q_aH},
\qquad
\omega_a=\frac{c_a^2}{q_aH}.
\]
Then
\[
\omega_a\ge0,quad \sum_a\omega_a=1,quad c^\top z=1,quad
\sum_aq_az_a^2=H^{-1},quad
z=\sum_a\omega_a\frac{e_a}{c_a}.
\tag{40}
\]
If \(B^{-1}z=(1,\beta_1,\ldots,\beta_{d-1})\), the convex-combination identity in (40) gives
\[
|\beta_j|\le F_j:=\max_{a\in S_c}
\left|\left(B^{-1}e_a/c_a\right)_{j+1}\right|.
\]
Thus \(D_q=\partial_x+\sum_j\beta_j\partial_{h_j}\) has
\[
J_{D_q}=J_x+J_h\sum_j\beta_j^2
\le J_{\max}:=J_x+J_h\sum_jF_j^2,
\tag{41}
\]
uniformly even near a simplex face.  Equations (35) and (40) imply \(\sum_aq_az_aC_a=\sigma^2/H\).  From (27),
\[
\left|\sum_aq_a\mathbb E(z_aC_a/V_a)-4\mathbb E\sigma^2/H\right|
\le64\alpha^3\Gamma L^2/H.
\tag{42}
\]
The normalization assumptions give the exact allocation identity
\[
\Delta:=H-4
=4\sum_{a\in S_c}\frac{(q_a-q_a^\star)^2}{q_a}+4(1-Q).
\tag{43}
\]
If (34) fails, (43) implies \(\Delta\ge\Delta_0:=4\delta^2\): an active deviation contributes at least \(4\delta^2\), while an inactive coordinate larger than \(\delta\) makes \(1-Q>\delta\ge\delta^2\).

Let \(W=\kappa(c)T+\varepsilon\sum_jk_j\), \(j_n=\alpha J_{\max}\), and \(D_0=4(1+u)+j_nH\).  By (23), \(0\le\sigma^2\le\alpha W\).  Insert (42) into (33), use
\[
\sum_aq_a\mathbb E(z_a^2/V_a)\le4(1+u)/H,
\]
and combine with (29).  Multiplication by \(nD_0\) and expansion of the square yield
\[
\begin{split}
D_0\{n\operatorname{BR}_{\Pi_n}(r)-1\}\ge{}&
\Delta-j_nH-4u-8\alpha W-128\alpha^3\Gamma L^2\\
&-32\alpha^2\Gamma^2(1+u)-8\alpha^2\Gamma^2j_nH.
\end{split}
\tag{44}
\]
For transparency, before the conservative bounds the leading expression is
\[
H(1-s)^2-D_0+D_0\left(\mathbb E\sigma^2-sum_aq_a\mathbb E(C_a^2/V_a)\right),
\]
where \(s=4\mathbb E\sigma^2/H+e\) and \(|He|\le64\alpha^3\Gamma L^2\).  Expanding this display, dropping the nonnegative square, using \(D_0-8\ge-8\), \(1+u\le2\), and (26)--(27) gives exactly (44).

The right side of (44) is nonnegative whenever
\[
8\alpha^2\Gamma^2\le1,
\qquad
j_n(1+4/\Delta_0)\le1/4,
\tag{45}
\]
and
\[
4u+8\alpha W+128\alpha^3\Gamma L^2
+32\alpha^2\Gamma^2(1+u)\le\Delta_0/2.
\tag{46}
\]
Indeed, \(H=4+\Delta\le\Delta(1+4/\Delta_0)\); (45) bounds the two terms involving \(j_nH\) by \(\Delta/2\), and (46) bounds all other negative terms by \(\Delta/2\).  These tests hold eventually because \(\alpha=n^{-1/3}\), \(\Delta_0=4n^{-1/4}\), and the slowest ratio is \(\alpha/\Delta_0=n^{-1/12}\).  Hence every positive exterior allocation satisfies
\[
\operatorname{BR}_{\Pi_n}(r)\ge1/n.
\tag{47}
\]

Finally, suppose \(q_a=0\) for some active arm \(a\).  Take \(z=e_a/c_a\) and write \(B^{-1}z=(1,\beta)\); the first coordinate is one because the first row of \(B^{-1}\) is \(c^\top\).  Every observed active marginal then has directional derivative zero, so the likelihood contribution in (32) is exactly zero and the numerator in (31) is \(\alpha\).  The finite coordinate bounds used in (41) give \(J_D\le J_{\max}\).  Consequently, (29) and (33) imply
\[
\operatorname{BR}_{\Pi_n}(r)\ge\alpha^2/J_{\max}\ge1/n
\tag{48}
\]
once \(\alpha J_{\max}\le1\).  This includes assignments supported entirely on inactive arms.  Equations (38), (47), and (48) cover every count vector, including all support-two faces.

To obtain the sharp coefficient, select functions \(f_m\) for which the left side of (24) decreases to \(\Lambda_c\), select \(\varepsilon_m\downarrow0\) so that \(\varepsilon_m\sum_jk_j\to0\), and, for each fixed \(m\), choose \(N_m\) beyond which (22), (39), and (45)--(48) hold with total coefficient error at most \(1/m\).  A diagonal choice \(m=m(n)\uparrow\infty\) sufficiently slowly gives a sequence of proper priors satisfying
\[
\inf_r\operatorname{BR}_{\Pi_n}(r)
\ge n^{-1}-\Lambda_cn^{-4/3}-o(n^{-4/3}).
\]

For any possibly dependent outcome-independent design, \(Z\) is independent of the prior schedule.  Conditional on a labeled assignment \(Z=z\), the preceding posterior calculation applies and permutation invariance makes its value depend only on \(r(z)\).  Under squared loss, conditional Jensen gives
\[
\mathbb E[(\widehat\tau-\tau_c)^2\mid O]
\ge\mathbb E[(\mathbb E[\widehat\tau\mid O]-\tau_c)^2\mid O]
\ge\operatorname{Var}(\tau_c\mid O),
\]
so auxiliary estimator randomization cannot improve upon the posterior mean.  Averaging the countwise inequality under any law of \(Z\) proves the asserted uniform all-design and all-estimator conclusion.
%%% FIELD statement-universal-lower
For every fixed \(K\ge2\), every \(c\in\mathcal C_K\), every outcome-independent possibly dependent design \(\mathcal D_n\), and every possibly randomized full-label estimator \(\widehat\tau\),
\[
\max_{y\in\{0,1\}^{n\times K}}\mathcal R_n(\mathcal D_n,\widehat\tau;y)
\ge n^{-1}-\Lambda_cn^{-4/3}+o(n^{-4/3}),
\]
where the remainder is uniform over the design-estimator pair. Equivalently, \(\limsup_{n\to\infty}n^{4/3}d_n(c)\le\Lambda_c\).  This includes dependent assignments supported on inactive-arm or missing-active-arm faces.
%%% FIELD proof-universal-lower
Let \(\Pi_n\) be the finite-schedule prior supplied by \(\mathrm{thm:all\text{-}count\text{-}bayes\text{-}inequality}\).  For any outcome-independent design \(\mathcal D_n\) and randomized estimator \(\widehat\tau\), the maximum risk dominates its prior average:
\[
\max_y\mathcal R_n(\mathcal D_n,\widehat\tau;y)
\ge\sum_y\Pi_n(y)\mathcal R_n(\mathcal D_n,\widehat\tau;y).
\tag{49}
\]
Outcome independence makes the assignment law independent of \(y\).  Condition on the entire assignment vector.  Conditional Jensen under squared loss removes auxiliary estimator randomization, and posterior optimality of the conditional mean then gives
\[
\sum_y\Pi_n(y)\mathcal R_n(\mathcal D_n,\widehat\tau;y)
\ge\sum_z\mathcal D_n(z)\operatorname{BR}_{\Pi_n}(r(z))
\ge\inf_r\operatorname{BR}_{\Pi_n}(r).
\tag{50}
\]
The countwise theorem covers every \(r(z)\), including inactive-only and missing-active-arm faces, and its remainder is independent of the design and estimator.  Combining (49)--(50) gives
\[
\max_y\mathcal R_n(\mathcal D_n,\widehat\tau;y)
\ge n^{-1}-\Lambda_cn^{-4/3}-o(n^{-4/3}).
\]
This is the theorem's displayed form because the negative of an \(o(n^{-4/3})\) sequence is again \(o(n^{-4/3})\).  Taking the infimum over the design-estimator pair proves the lower bound for \(R_n(c)\).  Since \(d_n(c)=n^{-1}-R_n(c)\), rearrangement yields \(\limsup_n n^{4/3}d_n(c)\le\Lambda_c\).  Nothing in this argument uses \(|S_c|\ge3\).
%%% FIELD statement-two-arm-reduction
For every fixed \(K\ge2\) and every \(c\in\mathcal C_K\) with \(|S_c|=2\), the two active coefficients are \(+1\) and \(-1\), up to exchange of their labels.  Consequently \(q^\star\) assigns probability \(1/2\) to each active arm and zero to every inactive arm, and \(V_c(s_c/2)=1\).  The independent upper construction ignores inactive arms, while the all-count lower construction applies with one transverse coordinate and covers allocation to inactive arms and every missing-active-arm face.  Hence
\[
R_n(c)=n^{-1}-C_An^{-4/3}+o(n^{-4/3}).
\]
When \(K=2\), this is exactly the expansion of Sudijono, Dobriban, and Tchetgen Tchetgen (2026), Theorem 3.1; arbitrary additional inactive labels do not change this expansion.
%%% FIELD proof-two-arm-reduction
Let \(S_c=\{a_1,a_2\}\).  The zero-sum condition gives \(c_{a_2}=-c_{a_1}\), and the \(\ell^1\)-normalization gives \(2|c_{a_1}|=2\).  Thus \(|c_{a_1}|=|c_{a_2}|=1\); exchanging labels and, if necessary, changing the common sign identifies the active vector with \((1,-1)\).  The definition of \(q^\star\) therefore assigns probability \(1/2\) to each active label and zero to every inactive label.

Suppose first that the active vector is \((1,-1)\).  Then \(v=(1/2,-1/2)\).  The response type \(t=(1,0)\) has \(b_t=v\) and \(\theta_t=1\), so the one-atom representation gives \(V_c(v)\le1\).  Conversely, every feasible representation obeys
\[
1=c^\top v=\sum_t\lambda_t\theta_t.
\]
Here \(\theta_t\in\{-1,0,1\}\).  Therefore
\[
\sum_t\lambda_t\theta_t^2
=\sum_t\lambda_t|\theta_t|
\ge\left|\sum_t\lambda_t\theta_t\right|=1.
\]
Thus \(V_c(v)=1\).  Exchanging the active labels or negating \(c\) maps response types bijectively, preserves the squared costs \(\theta_t^2\), and maps \(v\) accordingly, so the same conclusion holds for every support-two contrast.

The upper theorem applies because the product design uses only the two active labels; the value of every inactive potential outcome is irrelevant to its observations and its target.  The all-count theorem applies with \(d=2\), hence one transverse vector, and its cases (38), (47), and (48) explicitly include positive exterior allocations, allocations to inactive labels, and a zero count on either active label.  The universal lower theorem therefore applies to every dependent outcome-independent design and every randomized estimator.  Since \(\Lambda_c=C_AV_c(v)^{2/3}=C_A\), the upper and lower limits coincide and give the displayed expansion.  For \(K=2\), the cited node \(\mathrm{lem:published\text{-}two\text{-}arm\text{-}airy\text{-}law}\) gives the identical formula.  Conversely, adding inactive labels cannot improve the lower bound just proved and cannot worsen the explicit upper design, so it leaves the minimax value unchanged to the asserted order.
%%% FIELD statement-universal-airy-law
Theorem 1. For every fixed \(K\ge2\) and every zero-sum, \(\ell^1\)-normalized contrast \(c\in\mathcal C_K\),
\[
R_n(c)=n^{-1}-C_AV_c(s_c/2)^{2/3}n^{-4/3}+o(n^{-4/3}),
\qquad
\lim_{n\to\infty}n^{4/3}d_n(c)=\Lambda_c.
\]
The independent allocation \(q_a^\star=|c_a|/2\) on active arms and the total full-label nonlinear Airy estimator \(\delta_n^\star\) attain the expansion.  If \(|S_c|=2\), the reduction proposition gives \(V_c(s_c/2)=1\), allows any number of inactive labels, and recovers the published two-arm law.  The new response-alphabet and all-allocation-face argument supplies the closure for \(|S_c|\ge3\).
%%% FIELD proof-universal-airy-law
Fix \(K\ge2\) and \(c\in\mathcal C_K\).  The universal upper theorem gives
\[
\liminf_{n\to\infty}n^{4/3}d_n(c)\ge\Lambda_c,
\]
whereas the universal lower theorem gives
\[
\limsup_{n\to\infty}n^{4/3}d_n(c)\le\Lambda_c.
\]
The two limits coincide.  Using the definitions of \(d_n(c)\) and \(\Lambda_c\) yields
\[
R_n(c)=n^{-1}-\Lambda_cn^{-4/3}+o(n^{-4/3})
=n^{-1}-C_AV_c(s_c/2)^{2/3}n^{-4/3}+o(n^{-4/3}).
\]
The upper theorem proves attainment by the stated product allocation and total nonlinear estimator for every \(c\in\mathcal C_K\), including support two.  If \(|S_c|=2\), the reduction proposition proves that the two active coefficients are \(+1\) and \(-1\), that the gauge equals one, and that inactive labels neither improve the all-count lower bound nor obstruct the active-only upper design.  At \(K=2\) this agrees with the cited published theorem.  If \(|S_c|\ge3\), the same upper argument and the full-dimensional response-tube lower argument give the new response-alphabet closure.  Since zero sum and nonzero \(\ell^1\)-norm force \(|S_c|\ge2\), these cases exhaust the single natural class \(\mathcal C_K\).
%%% FIELD statement-effective-certificate
For every rational \(c\in\mathcal C_K\cap\mathbb Q^K\) and every rational \(\eta>0\), \(\mathsf{Cert}(c,\eta)\) terminates with a finite symbolic record \(\Gamma\). The record contains \(N\), finite rational and real-algebraic data, proof-checkable monotone bounds, and an explicitly specified total evaluator \(\mathsf{Eval}_{\Gamma}\). For every requested integer \(n\ge N\),
\[
\mathsf{Eval}_{\Gamma}(n)=(B_n,U_n)\in\mathbb A_{\mathrm R}^2,
\qquad B_n\le R_n(c)\le U_n,
\qquad n^{4/3}(U_n-B_n)\le\eta.
\]
Each \(B_n\) is the exact Bayes risk of a finite-schedule prior valid against all outcome-independent designs, and each \(U_n\) is the exact worst-schedule risk of a rationalized total full-label rule under the independent allocation \(q^\star\). Real-algebraic values are represented by square-free integer polynomials, rational isolating intervals, and Thom-sign data, and every arithmetic or order operation is certified by exact rational algorithms. As a continuity corollary, every real \(c\in\mathcal C_K\) admits normalization-preserving rational brackets \(c'\) with the same active support and signs such that
\[
|\sqrt{R_n(c)}-\sqrt{R_n(c')}|\le\omega(c,c').
\]
Choosing bracket errors \(o(n^{-5/6})\) transfers the coefficient and risk expansion mathematically to \(c\), without asserting a Turing evaluator for an unrepresented real contrast.  The certificate and transfer include support-two contrasts and arbitrary inactive arms.
%%% FIELD proof-effective-certificate
Fix rational \(c\in\mathcal C_K\) and rational \(\eta>0\).  We give a finite search, prove its termination, and identify the exact witnesses it returns.

The response-gauge lemma enumerates rational primal and dual basic feasible solutions and returns exact rational \(\kappa(c)>0\), \(\ell^\star\), and \(g^\star\).  The basis (20), nonnegative response representations of its columns, and the constants \(L_v,L_j,k_j,F_j,L,\Gamma\) are obtained using rational arithmetic and maxima over a finite response alphabet.  Thus all non-Airy constants in (21)--(48) are exact rational data.  When \(|S_c|=2\), the same enumeration has one transverse coordinate and the reduction proposition additionally certifies \(\kappa(c)=1\).

The one-dimensional Airy data have certified rational enclosures.  A concrete procedure uses the power-series recurrence for the decaying solution of \(A''(x)=xA(x)\), rational interval quadrature for its initial constants, and interval Taylor remainders on successive compact intervals.  Interval signs together with the differential equation isolate the rightmost negative zero of \(A'\); that zero is simple because there \(A''=xA\ne0\).  Bisection converges to the zero.  The same recurrence, followed by scaling with rational \(\kappa(c)\), gives arbitrarily narrow rational enclosures of \(C_A\), \(\Lambda_c\), \(h_\kappa\), \(h_\kappa'\), and valid tail and Lipschitz majorants.  Each requested enclosure is verified by finitely many rational inequalities, and refinement terminates by the series remainder bounds and isolating bisection.

Choose rational \(\epsilon_0<\eta/32\).  Enumerate compact rational piecewise-polynomial functions until exact integration and the Airy enclosure certify
\[
4\frac{\int f'^2}{\int f^2}
+\kappa(c)\frac{\int|x|f^2}{\int f^2}
\le\Lambda_c+\epsilon_0.
\tag{51}
\]
The approximation argument establishing (24) proves termination.  Choose rational \(\varepsilon>0\) satisfying \(\varepsilon\sum_jk_j\le\epsilon_0\).  Every tube probability is piecewise polynomial in \((x,h)\), with coefficients in \(\mathbb Q(n^{-1/3})\).  Exact integration over finitely many rational sign cells therefore makes every schedule mass a represented real-algebraic number.

For the upper rule, interval evaluation of \(h_\kappa\) at each algebraic score is refined until a rational approximation of the required accuracy is selected, and projection is then exact.  Let \(\widetilde\delta_n\) be the resulting finite rational table.  Because the target and both reported values lie in \([-1,1]\), an output perturbation of at most \(n^{-2}\) changes squared loss by at most \(4n^{-2}\).  The upper proof with its explicit Taylor bounds consequently supplies a proof-checkable threshold above which
\[
\max_y\mathcal R_n(\mathcal D_n^\star,\widetilde\delta_n;y)
\le n^{-1}-(\Lambda_c-4\epsilon_0)n^{-4/3}.
\tag{52}
\]
Likewise, (38), (47), (48), (51), and the choice of \(\varepsilon\) give a proof-checkable threshold above which
\[
\inf_r\operatorname{BR}_{\Pi_n}(r)
\ge n^{-1}-(\Lambda_c+4\epsilon_0)n^{-4/3}.
\tag{53}
\]
Every sufficient condition in (52)--(53) is a monotone inequality between explicit rational constants and negative rational powers of \(n\).  Sequential search therefore finds an integer \(N\) after which all of them hold.  Enlarge \(N\) to absorb the \(4n^{-2}\) rounding term and ensure \(8\epsilon_0\le\eta\).

We now define exact endpoints.  Let \(\nu_n(y)\) be the exact algebraic mass of schedule \(y\) under the selected tube prior.  For a fixed assignment \(z\), let \(O_z(y)\) be its observed string and set
\[
D_{z,o}=\sum_{y:O_z(y)=o}\nu_n(y).
\]
The exact conditional posterior Bayes risk is
\[
\operatorname{BR}_{\nu_n}(z)
=\sum_y\nu_n(y)\tau_c(y)^2
-\sum_{o:D_{z,o}>0}
\frac{\left\{\sum_{y:O_z(y)=o}\nu_n(y)\tau_c(y)\right\}^2}{D_{z,o}}.
\tag{54}
\]
Define
\[
B_n=\min_{z\in\mathcal A_K^n}\operatorname{BR}_{\nu_n}(z).
\tag{55}
\]
For any outcome-independent assignment law, posterior risk is the average of the conditional quantities in (54), so it is at least (55).  Hence \(B_n\le R_n(c)\).  Define
\[
U_n=\max_{y\in\{0,1\}^{n\times K}}
\sum_{z\in\mathcal A_K^n}\left(\prod_iq_{z_i}^\star\right)
\{\widetilde\delta_n(z,O_z(y))-\tau_c(y)\}^2.
\tag{56}
\]
This is the worst-schedule risk of one admissible design-estimator pair, so \(R_n(c)\le U_n\).  All sets in (54)--(56) are finite.  Field operations and order comparisons on algebraic entries terminate using resultants, isolating intervals, Thom signs, and Sturm sequences.  Hence \(B_n,U_n\in\mathbb A_{\mathrm R}\) are exact, and (52)--(53) imply, for every \(n\ge N\),
\[
B_n\le R_n(c)\le U_n,
\qquad
n^{4/3}(U_n-B_n)\le8\epsilon_0\le\eta.
\]
The finite constants, interval certificates, monotone inequalities, and algorithms (54)--(56) form \(\Gamma\) and specify a total evaluator on every integer \(n\ge N\).  This proves termination and exactness, without claiming a complexity bound.

It remains to justify the real-contrast transfer.  If \(c,c'\) have the same support and signs and satisfy the two normalizations, then \(u=c-c'\) has zero sum.  For every binary vector \(t\), writing \(P=\{a:u_a>0\}\) gives
\[
u^\top t\le\sum_{a\in P}u_a=\frac12\|u\|_1,
\]
and applying the same argument to \(-u\) yields
\[
|u^\top t|\le\frac12\|u\|_1=\omega(c,c').
\tag{57}
\]
For any fixed design-estimator pair, the triangle inequality in \(L^2\), (57), and the supremum over schedules show that the two root worst-case risks differ by at most \(\omega(c,c')\).  Applying this once to an arbitrarily near-optimal pair for \(c\) and once to an arbitrarily near-optimal pair for \(c'\), then letting the optimization errors vanish, proves
\[
|\sqrt{R_n(c)}-\sqrt{R_n(c')}|\le\omega(c,c').
\tag{58}
\]
Normalization-preserving rational contrasts with fixed support and signs are dense: the positive magnitudes lie in the relative interior of a rational simplex of mass one, and the negative magnitudes lie in another.  In the support-two case each simplex is a singleton, so the contrast is already rational up to its inactive zeros.

Local uniformity follows from finite geometry.  Inside a fixed sign chamber, the constraint matrix and right side of the gauge program are constant, and its value is the minimum, over finitely many basic feasible solutions \(\lambda\), of
\[
\sum_t\lambda_t(c^\top t)^2.
\]
Thus \(\kappa(c)\), and hence \(\Lambda_c\), is locally Lipschitz.  On a sufficiently small compact neighborhood of a fixed contrast, active \(|c_a|\) and \(\kappa(c)\) are bounded away from zero.  The finitely many primal masses, dual boxes, entries of \(B^{-1}\), response representations, and constants in the upper and lower remainders are then uniformly bounded.  The preceding upper and lower proofs therefore have a remainder \(o(n^{-4/3})\) uniformly on that neighborhood.

Choose rational \(c^{(n)}\) in the same chamber with \(e_n=\omega(c,c^{(n)})=o(n^{-5/6})\).  The unshrunk score estimator under \(q^\star\) has risk at most \(1/n\) by (7), (8), and (13), uniformly over the contrast.  Hence (58) gives
\[
|R_n(c)-R_n(c^{(n)})|
\le e_n\{2n^{-1/2}+e_n\}
=o(n^{-4/3}).
\tag{59}
\]
The locally uniform rational-contrast expansion and \(\Lambda_{c^{(n)}}\to\Lambda_c\) now transfer the coefficient and expansion to real \(c\).  Notice that this step uses the analytic theorem for the rational sequence, not a uniform running-time bound for its certificates.  It is a mathematical approximation argument only and does not assert a Turing evaluator for an unrepresented real number.
%%% FIELD statement-mixed-sign
For \(c=(1/3,2/3,-1/4,-3/4)\), exact primal and dual response-alphabet witnesses over all sixteen \(t\in\mathcal T_c\) give \(\kappa(c)=11/36\). The basis \(B=[v,d_1,d_2,d_3]\), with \(d_j\in\ker(c^\top)\), represents every missing-active-arm target direction, and adding an inactive fifth arm changes the allocation identity to
\[
H(q)-4=4\sum_{a\in S_c}(q_a-q_a^\star)^2/q_a+4\left(1-\sum_{a\in S_c}q_a\right).
\]
